% !TEX program = xelatex
% GMN3 中译 —— 主控导言区（自 GMN2 改写，保留原文字母宏以便逐式核对）
\usepackage[a4paper,margin=2.6cm,top=2.8cm,bottom=2.8cm]{geometry}
\usepackage{amsmath,amssymb,amsfonts,amsthm,mathtools}
\usepackage{mathrsfs}
\usepackage{bm}
\usepackage{graphicx}
\usepackage{booktabs}
\usepackage{array}
\usepackage{enumitem}
\usepackage{xcolor}
\usepackage{tcolorbox}
\usepackage{hyperref}
\usepackage{xspace}
\usepackage{tikz}
\usepackage{epigraph}

\graphicspath{{figures/}}

% ---------- 字体与行距 ----------
\linespread{1.18}
\setlength{\parskip}{0.25em}
\emergencystretch=1.2em
\sloppy

% 公式按节编号（与原刊排版一致：(1.1)、(2.18)…）
\numberwithin{equation}{section}

% ---------- 超链接 ----------
\hypersetup{
  colorlinks=true,
  linkcolor=blue!45!black,
  citecolor=green!40!black,
  urlcolor=magenta!60!black,
  pdftitle={带框 BPS 态（Framed BPS States）},
  pdfauthor={Gaiotto--Moore--Neitzke},
}

% ---------- 定理类环境（与数译体例一致：各类独立连续编号） ----------
\theoremstyle{plain}
\newtheorem{thm}{定理}
\newtheorem{lem}[thm]{引理}
\newtheorem{prop}[thm]{命题}
\newtheorem{cor}[thm]{推论}
\newtheorem{conj}[thm]{猜想}
\theoremstyle{definition}
\newtheorem{defn}{定义}
\newtheorem{exmp}{例}
\newtheorem{prob}{问题}
\theoremstyle{remark}
\newtheorem{rem}{注}
\newtheorem*{conv}{约定}
\newtheorem*{notn}{记号}

% ---------- 译注 ----------
\newcounter{transnote}
\newenvironment{yizhu}
  {\par\medskip\refstepcounter{transnote}%
   \begin{tcolorbox}[colback=orange!4,colframe=orange!55!black,boxrule=0.5pt,
                     left=4pt,right=4pt,top=3pt,bottom=3pt,arc=1pt]
   \noindent\textbf{译注 \thetransnote：}}
  {\end{tcolorbox}\par\medskip}
% 行内短译注：\zy{...} 生成脚注
\newcommand{\zy}[1]{\footnote{\textcolor{orange!70!black}{【译注】}#1}}

% ---------- 原文字母宏（与 wcf-paper.tex 对齐，便于逐式核对） ----------
\newcommand{\IC}{\mathbb{C}}
\newcommand{\IN}{\mathbb{N}}
\newcommand{\IZ}{\mathbb{Z}}
\newcommand{\IR}{\mathbb{R}}
\newcommand{\IP}{\mathbb{P}}
\newcommand{\IQ}{\mathbb{Q}}

\newcommand{\CA}{\mathcal{A}}
\newcommand{\CB}{\mathcal{B}}
\newcommand{\CC}{\mathcal{C}}
\newcommand{\CD}{\mathcal{D}}
\newcommand{\CE}{\mathcal{E}}
\newcommand{\CF}{\mathcal{F}}
\newcommand{\CG}{\mathcal{G}}
\newcommand{\CH}{\mathcal{H}}
\newcommand{\CI}{\mathcal{I}}
\newcommand{\CJ}{\mathcal{J}}
\newcommand{\CK}{\mathcal{K}}
\newcommand{\CL}{\mathcal{L}}
\newcommand{\CM}{\mathcal{M}}
\newcommand{\CN}{\mathcal{N}}
\newcommand{\CO}{\mathcal{O}}
\newcommand{\CP}{\mathcal{P}}
\newcommand{\CQ}{\mathcal{Q}}
\newcommand{\CR}{\mathcal{R}}
\newcommand{\CS}{\mathcal{S}}
\newcommand{\CT}{\mathcal{T}}
\newcommand{\CU}{\mathcal{U}}
\newcommand{\CV}{\mathcal{V}}
\newcommand{\CW}{\mathcal{W}}
\newcommand{\CX}{\mathcal{X}}
\newcommand{\CY}{\mathcal{Y}}
\newcommand{\CZ}{\mathcal{Z}}

\newcommand{\N}{\mathcal{N}}
\newcommand{\la}{\langle}
\newcommand{\ra}{\rangle}
\newcommand{\half}{\tfrac{1}{2}}
\newcommand{\abs}[1]{\lvert #1 \rvert}
\newcommand{\norm}[1]{\lVert #1 \rVert}
\newcommand{\inprod}[1]{\langle #1 \rangle}
\newcommand{\pz}{(\zeta)}
\newcommand{\ti}[1]{\textit{#1}}
\newcommand{\re}{\operatorname{Re}}
\newcommand{\im}{\operatorname{Im}}
\newcommand{\tr}{\operatorname{tr}}
\newcommand{\sgn}{\operatorname{sgn}}
\newcommand{\vol}{\operatorname{vol}}
\newcommand{\Res}{\operatorname{Res}}
\newcommand{\Tr}{\operatorname{Tr}}
\newcommand{\sfm}{\mathrm{sf}}
\newcommand{\hk}{超凯勒}
\newcommand{\qk}{四元数凯勒}
\newcommand{\kahler}{凯勒}
\newcommand{\elec}{\mathrm{e}}
\newcommand{\magn}{\mathrm{m}}
\newcommand{\degen}{\Omega}
\newcommand{\hsymp}{\varpi}
\newcommand{\clat}{\Gamma}
\newcommand{\vm}{\CB}
\newcommand{\p}{\partial}
\newcommand{\pb}{\bar\partial}
\newcommand{\inst}{\mathrm{inst}}
\newcommand{\elecm}{\mathrm{e}}
\newcommand{\cwarrow}{\text{\Large$\curvearrowright$}}
\newcommand{\ccwarrow}{\text{\Large$\curvearrowleft$}}

% —— 原文自定义宏（与原文同名，便于逐式核对；勿在各 chunk 里重复定义）——
\renewcommand\sf{{\mathrm{sf}}}
\def\tp{\tilde\phi}
\def\n{\vec{n}}
\def\nd{\vec{n}^{\dagger}}
\newcommand{\plog}[1]{#1\partial_{#1}}
\newcommand{\eps}{\epsilon}

% —— GMN3（1006.0146）追加的原文宏 ——
% 注意：原文另有 \def\be{\begin{equation}} / \def\ee{\end{equation}} 与大量单字母宏
% （\a \b \c \p \la \ra \half ...），一律不抄进来，避免覆盖译者常用命令。
\newcommand{\RH}{{\rm RH}}
\newcommand{\WKB}{{\rm WKB}}
\newcommand{\hyper}{{\rm hyper}}
\newcommand{\vect}{{\rm vector}}
\newcommand{\core}{{\rm core}}
\newcommand{\halo}{{\rm halo}}
\newcommand{\BPS}{{\rm BPS}}
\newcommand{\fro}{\overline{\underline{\Omega}}}
\newcommand{\reg}{{\rm reg}}
\newcommand{\flav}{{\rm flavor}}
\newcommand{\juggle}{juggle\xspace}
\newcommand{\juggles}{juggles\xspace}
\newcommand{\juggleing}{juggling\xspace}
\newcommand{\nsing}{{l}}
\newcommand{\ts}{\tilde{s}}
\newcommand{\id}{\mathbf{1}}
\newcommand{\hkq}{/\!/\!/}
\newcommand{\bS}{\mathbf S}
\newcommand{\IS}{{\mathbb S}}
\newcommand{\tCM}{\widetilde{\cal M}}
\newcommand{\tCY}{\widetilde{\cal Y}}
\newcommand{\lieg}{\mathfrak{g}}
\newcommand{\liet}{\mathfrak{t}}
% 原文 \sect{...}：重设公式计数器后再起节（原文用它实现按节编号）
\newcommand{\sect}[1]{\setcounter{equation}{0}\section{#1}}

\DeclareMathOperator{\rank}{rank}
\DeclareMathOperator{\Area}{Area}
\DeclareMathOperator{\End}{End}
\DeclareMathOperator{\Pexp}{Pexp}
\DeclareMathOperator{\sTr}{sTr}
\DeclareMathOperator{\Hom}{Hom}
\DeclareMathOperator{\Hol}{Hol}
\DeclareMathOperator{\sdim}{sdim}
\DeclareMathOperator{\Edges}{Edges}
\DeclareMathOperator{\fund}{fund}

% 插图：原文 \insfig{名}{题注} 对应 figures/名-crop.pdf（scale=0.275）。
% 本译本改用 \insfig[宽]{名}{题注}，默认 0.8\textwidth，便于统一版面。
\newcommand{\insfig}[3][0.8\textwidth]{\begin{figure}[htbp]\centering\includegraphics[width=#1]{figures/#2-crop.pdf}\caption{#3}\label{fig:#2}\end{figure}}

% 原文 \instwofigs{宽1}{名1}{注1}{宽2}{名2}{注2}
\newcommand{\instwofigs}[6]{
\begin{figure}[htbp]
\begin{minipage}[b]{0.5\linewidth}
\centering \includegraphics[width=#1]{figures/#2-crop.pdf} \caption{#3} \label{fig:#2}
\end{minipage}
\hspace{0.4cm}
\begin{minipage}[b]{0.45\linewidth}
\centering \includegraphics[width=#4]{figures/#5-crop.pdf} \caption{#6} \label{fig:#5}
\end{minipage}
\end{figure}
}

% 原文 \insthreefigs{高1}{名1}{注1}{高2}{名2}{注2}{高3}{名3}{注3}
\newcommand{\insthreefigs}[9]{
\begin{figure}[htbp]
\begin{minipage}[b]{0.31\linewidth}
\centering \includegraphics[height=#1]{figures/#2-crop.pdf} \caption{#3} \label{fig:#2}
\end{minipage}
\hspace{0.25cm}
\begin{minipage}[b]{0.31\linewidth}
\centering \includegraphics[height=#4]{figures/#5-crop.pdf} \caption{#6} \label{fig:#5}
\end{minipage}
\hspace{0.25cm}
\begin{minipage}[b]{0.31\linewidth}
\centering \includegraphics[height=#7]{figures/#8-crop.pdf} \caption{#9} \label{fig:#8}
\end{minipage}
\end{figure}
}

% 译者常用的几家译名（统一在此，便于全系列一致）
\newcommand{\KS}{Kontsevich--Soibelman\xspace}
\newcommand{\SW}{Seiberg--Witten\xspace}
\newcommand{\CFT}{CFT\xspace}

% 原文自定义宏（ch04/ch05 用到，原稿未定义）
\newcommand{\ch}{{\rm ch}}
\providecommand{\hW}{\widehat W}
\newcommand{\qbinom}[2]{\genfrac{[}{]}{0pt}{}{#1}{#2}}
